\section{Research Objectives}

The overarching goal of this study is to design and analyse a detector-agnostic
framework for multi-person pose tracking that couples a depth-aware graph
representation with a Graph Attention Network (GAT).  
Within this broad aim, the work pursues four concrete objectives:

\begin{itemize}
    \item \textbf{Depth-Aware Graph Construction}\,:  
    Formulate a graph in which each node stores a joint’s image coordinates and
    estimated depth while each edge encodes spatial and geometric relations;
    analyse how the added depth channel alters feature compactness and
    expressiveness compared with purely 2-D graphs.

    \item \textbf{Accurate Joint-to-Person Assignment}\,:  
    Investigate whether incorporating depth into the GAT improves joint
    grouping accuracy in crowded scenes where limbs from different individuals
    overlap or appear in close proximity.

    \item \textbf{Computational Efficiency}\,:  
    Evaluate whether the compact, graph-based feature space shortens training
    time and sustains real-time inference, identifying trade-offs between model
    depth, runtime, and memory footprint.

    \item \textbf{Detector-Agnostic Scalability}\,:  
    Demonstrate that the framework can ingest keypoints produced by diverse
    detectors—ranging from lightweight CNN heads to transformer
    architectures—without architectural changes, and analyse performance
    sensitivity to the underlying detector.
\end{itemize}

Together, these objectives target a pose-tracking solution that is accurate in
crowded settings, computationally practical, and readily transferable across 
pose detection front-ends.